\documentclass[11pt]{article}

\usepackage[utf8]{inputenc}
\usepackage[T1]{fontenc}
\usepackage{lmodern}
\usepackage{geometry}
\usepackage{amsmath,amssymb,amsfonts,amsthm,mathtools}
\usepackage{array}
\usepackage{enumitem}
\usepackage{microtype}
\usepackage{hyperref}

\geometry{margin=1in}
\hypersetup{colorlinks=true, linkcolor=black, urlcolor=blue, citecolor=black}
\setlist[enumerate]{topsep=0.35em,itemsep=0.3em}
\setlist[itemize]{topsep=0.3em,itemsep=0.25em}

\newcommand{\Lang}{\mathcal{L}_{0}}
\newcommand{\Obj}{\mathsf{Obj}_{0}}
\newcommand{\Inv}{\mathsf{Inv}_{0}}
\newcommand{\Cmp}{\mathsf{Cmp}_{0}}
\newcommand{\EqSrc}{\mathsf{Eq}_{\mathrm{src}}}
\newcommand{\GenH}{\mathsf{Gen}_{H}}
\newcommand{\RetainH}{\mathsf{Retain}_{H}}
\newcommand{\NoTargetH}{\mathsf{NoTarget}_{H}}
\newcommand{\NoTargetSyn}{\mathsf{NoTarget}_{H}^{\mathrm{syn}}}
\newcommand{\SrcTok}{\mathsf{Tok}_{0}}
\newcommand{\ForbidT}{\mathsf{Forbid}_{T}}
\newcommand{\ProxyT}{\mathsf{Proxy}_{T}}
\newcommand{\TraceH}{\mathsf{Trace}_{H}}
\newcommand{\Annot}{\mathsf{Ann}}
\newcommand{\Influence}{\mathsf{Infl}}
\newcommand{\FreezeH}{\mathsf{Freeze}_{H}}
\newcommand{\JudgeSyn}{\mathsf{J}_{\mathrm{syn}}}
\newcommand{\Pass}{\mathsf{Pass}}
\newcommand{\Block}{\mathsf{Block}}
\newcommand{\Pause}{\mathsf{Pause}}

\newtheorem{definitionblock}{Definition}
\newtheorem{obligation}{Obligation}
\newtheorem{finding}{Finding}
\newtheorem{guard}{Guard}
\newtheorem{verdict}{Verdict}
\newtheorem{nextburden}{Next burden}

\title{\textbf{Localization Source-Basis Axiom Selector-Domain EqSrc Source-Equivalence Obligation Packet or Controlled Pause}\\[0.35em]
\large Ontology Formalizer Review After the Stress-Tested \(\NoTargetSyn\) Source-Syntax Guard}
\author{Ontology Formalizer AgentJob AJ-RT-20260613-025-001}
\date{June 14, 2026}

\begin{document}

\maketitle

\begin{abstract}
This draft/control artifact states the next \(\EqSrc\)
source-equivalence obligation after the stress-tested
\(\NoTargetSyn\) source-syntax guard. The prior local guard is useful:
it blocks target-channel imports, unmapped proxies, weakened annotations,
implicit influence-edge absence, and semantic target comparison. It is not,
however, a source-equivalence relation. The available registered surfaces do
not yet supply a source-only domain of compared objects, a witness-pair
ledger, source-language invariants, a comparison rule, and step-level
proof that source-equivalence decisions are not selected by target proxies.
The result is therefore a controlled pause with a precise obligation packet.
It does not discharge \(\EqSrc\), \(\RetainH\), or \(\GenH\), and does not
authorize candidate reconstruction, canonical ontology edit, benchmark
promotion, Gate Chair review, or completed-derivation language.
\end{abstract}

\section{Scope}

This artifact is the role output for task \texttt{RT-20260613-025}. It
answers the control question requested by \texttt{handoff-0041}: what must
exist before \(\EqSrc\) can be treated as a source-equivalence result?

It does not edit canonical ontology TeX, does not introduce permanent
vocabulary into \(\Lang\), does not prove semantic \(\NoTargetH\), and does
not convert the local \(\NoTargetSyn\) source-syntax \(\Pass\) into
\(\EqSrc\) discharge.

\section{Available Input}

\begin{definitionblock}[Stress-tested guard]
The prior Refuter result supports only the local statement that the
manifest-bound \(\NoTargetSyn\) source-syntax \(\Pass\) survives
benchmark-label erasure and structure-preserving source-neutral relabeling,
while proxy omissions, new target-success proxies, influence perturbations,
annotation weakening, and semantic target comparison return \(\Block\) or
require a superseding manifest.
\end{definitionblock}

\begin{guard}[Guard input is not equivalence discharge]
The guard may be used to reject target-channel leakage in a future
\(\EqSrc\) attempt. It may not be used as the equivalence relation itself.
Syntactic no-target admissibility and source-equivalence are distinct
burdens: the first blocks inadmissible channels; the second must compare
source-side objects by source-only invariants.
\end{guard}

\section{EqSrc Obligation Packet}

\begin{definitionblock}[Minimal source-equivalence datum]
A future \(\EqSrc\) attempt requires a registered source-only datum
\[
    \mathcal{E}_{0} = (\Obj,\Inv,\Cmp,\Annot,\Influence),
\]
where:
\begin{enumerate}
    \item \(\Obj\) is a finite declared domain of source-language objects or
    object families being compared.
    \item \(\Inv\) is a finite ledger of source-side invariants used by the
    comparison.
    \item \(\Cmp\) is a source-only comparison rule whose premises do not
    use target metrics, target proper time, target locality, target
    observer protocols, target-favorable rank, benchmark-success evidence,
    or canonical adoption labels.
    \item \(\Annot\) binds every comparison step to source path, object
    identifier, hash, timing, rule scope, dependency, and proxy-exclusion
    content.
    \item \(\Influence\) records the influence edges by which forbidden
    channels and proxy classes are blocked rather than silently ignored.
\end{enumerate}
\end{definitionblock}

\begin{obligation}[Domain and witness pairs]
The packet must name at least one source-side witness pair
\[
    (x,y) \in \Obj \times \Obj
\]
and state whether the intended judgment is \(\EqSrc(x,y)\),
\(\neg\EqSrc(x,y)\), or \(\Block\). The pair must be selected before any
target-success criterion can influence the comparison.
\end{obligation}

\begin{obligation}[Invariant ledger]
For every compared witness pair, the packet must identify the source-only
invariants \(I_{1},\dots,I_{n} \in \Inv\) used by the comparison and state
why each invariant is source-side. A notation-similarity claim, validation
passage, rank, family selection, generated-derivative visibility, or human
authorization token is a proxy surface, not an invariant by itself.
\end{obligation}

\begin{obligation}[Comparison rule]
The comparison rule \(\Cmp(x,y;\Inv)\) must be mechanically checkable from
registered source artifacts. If the rule asks whether \(x\) and \(y\) have
the same target metric behavior, proper-time behavior, target locality,
observer protocol, benchmark rank, or target-favorable empirical status,
the result is \(\Block\), not \(\EqSrc\).
\end{obligation}

\begin{obligation}[Proxy and influence coverage]
The existing manifest identifies source-equivalence proxy PT-002, trace
event ST-002, annotation obligation AN-002, and influence edge IE-006. In a
future \(\EqSrc\) attempt these are not discharge evidence. They are the
minimum guard surfaces that must be instantiated with a comparison rule and
proxy-exclusion proof for each witness pair.
\end{obligation}

\begin{obligation}[Negative controls]
A future packet must include source-side negative controls: at least one
pair or class for which the same source comparison rule returns
\(\neg\EqSrc\) or \(\Block\) for source-side reasons. Without this control,
the rule risks collapsing into notation similarity, status agreement, or a
target-success proxy.
\end{obligation}

\section{Assessment of Current Evidence}

\begin{finding}[The current manifest contains obligation labels, not
equivalence witnesses]
The registered freeze manifest includes ST-002 as a source-equivalence
comparison event, PT-002 as a source-equivalence proxy, AN-002 as the
required comparison annotation, and IE-006 as the influence edge from PT-002
to ST-002. These entries specify where an \(\EqSrc\) attempt must be
checked. They do not provide a witness pair, invariant set, or comparison
rule.
\end{finding}

\begin{finding}[The source-syntax pass does not compare objects]
The \(\NoTargetSyn\) source-syntax \(\Pass\) from \texttt{RT-20260613-023}
and its stress test in \texttt{RT-20260613-024} evaluate whether a declared
trace and manifest avoid forbidden target channels and unmapped proxy use.
They do not compare two source objects under a source-only equivalence
relation.
\end{finding}

\begin{finding}[Source-neutral relabeling is necessary but insufficient]
The stress-test survival under structure-preserving relabeling is useful
evidence that token names are not doing the work. It is not sufficient for
\(\EqSrc\), because equivalence requires a specified comparison domain,
source invariants, and a rule that can be applied to witness pairs.
\end{finding}

\section{Controlled Pause}

\begin{verdict}[EqSrc controlled pause]
The available tracked source material does not yet support a completed
\(\EqSrc\) source-equivalence judgment. The correct result is
\[
    \EqSrc \Rightarrow \Pause
\]
with the obligation packet above preserved for a future source-equivalence
evidence-supply task.
\end{verdict}

\begin{verdict}[No downstream discharge follows]
This controlled pause does not discharge \(\EqSrc\), \(\RetainH\), or
\(\GenH\). It does not prove semantic \(\NoTargetH\), does not canonically
adopt \(\NoTargetSyn\), and does not authorize candidate reconstruction,
canonical ontology edit, benchmark promotion, Gate Chair review, or
completed-derivation language.
\end{verdict}

\section{Boundary Ledger}

\begin{center}
\begin{tabular}{p{0.24\linewidth}p{0.44\linewidth}p{0.18\linewidth}}
\textbf{Surface} & \textbf{Result} & \textbf{Status} \\
\hline
\(\NoTargetSyn\) source-syntax guard & usable only as blocker and proxy-leakage guard & retained \\
ST-002 comparison event & names an obligation surface & not discharge \\
PT-002 source-equivalence proxy & identifies leakage risk & blocked until instantiated \\
AN-002 comparison annotation & requires source-only comparison proof & missing \\
IE-006 influence edge & maps PT-002 to ST-002 & guard only \\
\(\Obj,\Inv,\Cmp\) & not registered as a complete source-equivalence datum & missing \\
\(\EqSrc\) & obligation packet stated & \(\Pause\) \\
\(\RetainH,\GenH\) & not attempted & blocked \\
Candidate reconstruction & not authorized & blocked \\
Gate Chair review & not authorized and human-gated & blocked \\
\end{tabular}
\end{center}

\section{Next Burden}

\begin{nextburden}[Evidence-supply before equivalence claim]
The logical next step is not Candidate Constructor or Gate Chair review. A
future continuation may open a bounded source-equivalence evidence-supply
packet only if it can register \(\Obj\), witness pairs, \(\Inv\), \(\Cmp\),
\(\Annot\), and \(\Influence\) without target-channel or proxy selection. If
that evidence cannot be supplied, the controlled pause should remain in
force. \(\RetainH\), \(\GenH\), candidate reconstruction, canonical ontology
edit, benchmark promotion, Gate Chair review, and completed-derivation
language remain blocked.
\end{nextburden}

\section*{References}

Ricciardi, A. S. (2026, June 14). \emph{Localization source-basis axiom
selector-domain NoTarget source-syntax benchmark-independence and
proxy-variation stress test: Refuter review of the manifest-bound
\(\NoTargetSyn\) source-syntax pass} [Unpublished manuscript]. The
Aether-Flow Interpretation of Relativity Research Project,
\texttt{research\_control/tasks/RT-20260613-024/artifacts/40\_LOCALIZATION\_SOURCE\_BASIS\_AXIOM\_SELECTOR\_DOMAIN\_NO\_TARGET\_SOURCE\_SYNTAX\_BENCHMARK\_INDEPENDENCE\_PROXY\_VARIATION\_STRESS\_TEST.tex}.

Ricciardi, A. S. (2026, June 13). \emph{Localization source-basis axiom
selector-domain NoTarget source-syntax theorem retry: Manifest-bound
\(\NoTargetSyn\) adjudication after source-registered freeze audit}
[Unpublished manuscript]. The Aether-Flow Interpretation of Relativity
Research Project,
\texttt{research\_control/tasks/RT-20260613-023/artifacts/39\_LOCALIZATION\_SOURCE\_BASIS\_AXIOM\_SELECTOR\_DOMAIN\_NO\_TARGET\_SOURCE\_SYNTAX\_THEOREM\_RETRY.tex}.

Ricciardi, A. S. (2026, June 13). \emph{Localization source-basis axiom
selector-domain NoTarget source-registered freeze manifest: Complete
draft/control inventory for Tok0, ForbidT, ProxyT, Ann, and Infl}
[Unpublished manuscript]. The Aether-Flow Interpretation of Relativity
Research Project,
\texttt{research\_control/tasks/RT-20260613-021/artifacts/37\_LOCALIZATION\_SOURCE\_BASIS\_AXIOM\_SELECTOR\_DOMAIN\_NO\_TARGET\_SOURCE\_REGISTERED\_FREEZE\_MANIFEST.tex}.

\end{document}
